\documentclass[10pt,a4paper,reqno]{amsart}

\usepackage[a4paper,margin=25mm]{geometry}
\usepackage{amsmath,amssymb,mathtools}
\usepackage{booktabs,array,tabularx}
\usepackage{enumitem,etoolbox,microtype}
\usepackage[hidelinks]{hyperref}
\usepackage{xurl}
\urlstyle{tt}
\allowdisplaybreaks[2]
\emergencystretch=3em
\setlist{nosep}
\AtBeginEnvironment{verbatim}{\footnotesize}

\newtheorem{theorem}{Theorem}[section]
\newtheorem{lemma}[theorem]{Lemma}
\newtheorem{proposition}[theorem]{Proposition}
\newtheorem{corollary}[theorem]{Corollary}
\theoremstyle{definition}
\newtheorem{definition}[theorem]{Definition}
\newtheorem{remark}[theorem]{Remark}

\newcommand{\C}{\mathbb C}
\newcommand{\Q}{\mathbb Q}
\newcommand{\F}{\mathbb F}
\newcommand{\Sym}{\operatorname{Sym}}
\newcommand{\rank}{\operatorname{rank}}
\newcommand{\im}{\operatorname{im}}
\newcommand{\coker}{\operatorname{coker}}
\newcommand{\Span}{\operatorname{span}}
\newcommand{\ChowRank}{\operatorname{ChowRank}}
\newcommand{\perm}{\operatorname{perm}}
\newcommand{\row}{\operatorname{row}}
\newcommand{\id}{\operatorname{id}}
\newcolumntype{L}{>{\raggedright\arraybackslash}X}

\title{The Chow Rank of the Three-by-Three Permanent}
\author{}
\date{September 1, 2026}
\subjclass[2020]{Primary 15A69; Secondary 14N05, 68V05}
\keywords{permanent, Chow rank, Koszul flattening}
\hypersetup{pdftitle={The Chow Rank of the Three-by-Three Permanent}}

\begin{document}
\begin{abstract}
We prove over every field of characteristic zero that
\[
  \ChowRank(\perm_3)=4.
\]
The lower bound follows from the kernel of the catalecticant--Koszul
flattening, and the upper bound is Glynn's four-term decomposition. The
argument is self-contained, and an independent exact reconstruction program
is included with the paper.
\end{abstract}
\maketitle

\section{The three-by-three permanent}
\subsection{Definitions and statement}

Let
\[
  V=\Span_{\C}\{x_{ij}:0\leq i,j\leq2\},\qquad \dim V=9,
\]
and define
\[
  \perm_3=\sum_{\sigma\in S_3}
  x_{0,\sigma(0)}x_{1,\sigma(1)}x_{2,\sigma(2)}\in\Sym^3V.
\]

\begin{definition}
The Chow rank of a cubic form \(f\in\Sym^3V\) is the least integer \(r\)
for which
\[
  f=\sum_{s=1}^{r}\ell_{s1}\ell_{s2}\ell_{s3},
  \qquad \ell_{st}\in V.
\]
\end{definition}

\begin{theorem}\label{n3:thm:main}
\[
  \boxed{\ChowRank(\perm_3)=4}.
\]
\end{theorem}

\subsection{The catalecticant and Koszul flattening}

For a cubic form \(f\in\Sym^3V\), define
\[
  C_{1,2}(f):V^*\longrightarrow\Sym^2V,
  \qquad D\longmapsto D(f),
\]
and write \(E_f:=\im C_{1,2}(f)\). Define
\[
  \delta:\Sym^2V\otimes V\longrightarrow V\otimes\bigwedge\nolimits^2V
\]
by
\begin{equation}\label{n3:eq:delta}
  \delta((uv)\otimes w)
  =v\otimes(u\wedge w)+u\otimes(v\wedge w),
\end{equation}
extended linearly. The associated Koszul flattening is
\[
  K(f):=\delta\circ\bigl(C_{1,2}(f)\otimes\id_V\bigr),
\]
so
\begin{equation}\label{n3:eq:image}
  \im K(f)=\delta(E_f\otimes V).
\end{equation}

\subsubsection{The Koszul kernel and first prolongation}

For \(E\subseteq\Sym^2V\), define
\[
  E^{(1)}:=\left\{g\in\Sym^3V:
  \frac{\partial g}{\partial x_{ij}}\in E
  \text{ for every }i,j\right\}.
\]

\begin{lemma}\label{n3:lem:kernel-prolongation}
There is a natural isomorphism
\[
  \ker\bigl(\delta|_{E\otimes V}\bigr)\simeq E^{(1)}.
\]
Consequently,
\[
  \rank\delta(E\otimes V)=9\dim E-\dim E^{(1)}.
\]
\end{lemma}

\begin{proof}
Define
\[
  \iota:\Sym^3V\longrightarrow\Sym^2V\otimes V,
  \qquad
  \iota(g)=\frac13\sum_{i,j}
  \frac{\partial g}{\partial x_{ij}}\otimes x_{ij}.
\]
Euler's identity shows that tensor multiplication sends \(\iota(g)\) to
\(g\), so \(\iota\) is injective. Applying \eqref{n3:eq:delta} to a cubic
monomial shows that the wedge terms cancel in pairs; hence
\(\delta\circ\iota=0\).

Conversely, write the coefficients of \(z\in\Sym^2V\otimes V\) as
\(z_{abc}\). Symmetry in the first two indices gives
\(z_{abc}=z_{bac}\), and \(\delta(z)=0\) gives
\(z_{abc}=z_{acb}\). These transpositions generate every permutation of the
three indices, so \(z_{abc}\) is fully symmetric. Thus \(z=\iota(g)\) for a
unique \(g\in\Sym^3V\), and
\[
  \ker\delta=\iota(\Sym^3V).
\]
Intersecting with \(E\otimes V\) says precisely that all first derivatives of
\(g\) lie in \(E\), which gives \(E^{(1)}\). Rank--nullity and \(\dim V=9\)
give the formula.
\end{proof}

\subsection{The flattening rank of the permanent}

Write \(E:=E_{\perm_3}\). Differentiation with respect to \(x_{ij}\)
produces the \(2\times2\) subpermanent obtained by deleting row \(i\) and
column \(j\). Therefore
\begin{equation}\label{n3:eq:E}
  E=\Span\left\{x_{ia}x_{jb}+x_{ib}x_{ja}:
  0\leq i<j\leq2,\ 0\leq a<b\leq2\right\}.
\end{equation}
Distinct row-pair and column-pair supports are disjoint, so these nine
quadrics are linearly independent:
\begin{equation}\label{n3:eq:E-dim}
  \dim E=\binom32^2=9.
\end{equation}

\begin{proposition}\label{n3:prop:prolongation}
\[
  E^{(1)}=\Span(\perm_3).
\]
In particular, \(\dim E^{(1)}=1\).
\end{proposition}

\begin{proof}
Every first derivative of \(\perm_3\) lies in \(E\), so
\(\perm_3\in E^{(1)}\).

Conversely, let \(g\in E^{(1)}\). Every quadratic monomial in
\eqref{n3:eq:E} uses distinct variable rows and columns. Thus
\(\partial g/\partial x_{ij}\in E\) forces every nonzero cubic monomial of
\(g\) to use three distinct rows and columns. Hence
\[
  g=\sum_{\sigma\in S_3}c_{\sigma}
  x_{0,\sigma(0)}x_{1,\sigma(1)}x_{2,\sigma(2)}.
\]
Fix one variable and differentiate with respect to it. On the remaining two
rows and columns there are two matchings. The corresponding part of \(E\) is
one-dimensional and generated by their sum, so the two coefficients are
equal. Thus \(c_\sigma=c_\tau\) whenever \(\sigma\) and \(\tau\) differ by a
transposition. Transpositions generate \(S_3\), so all coefficients are equal
and \(g\) is a scalar multiple of \(\perm_3\).
\end{proof}

\begin{corollary}\label{n3:cor:perm-rank}
\[
  \rank K(\perm_3)=80.
\]
\end{corollary}

\begin{proof}
Lemma~\ref{n3:lem:kernel-prolongation}, \eqref{n3:eq:E-dim}, and
Proposition~\ref{n3:prop:prolongation} give
\[
  \rank K(\perm_3)=9\cdot9-1=80.
\]
\end{proof}

\subsection{A bound for one Chow term}

\begin{proposition}\label{n3:prop:term-rank}
For every cubic Chow term \(T=\ell_1\ell_2\ell_3\),
\[
  \rank K(T)\leq26.
\]
Equality holds when the three factors are linearly independent.
\end{proposition}

\begin{proof}
Let \(F:=E_T\). Differentiating \(T\) gives
\[
  F\subseteq\Span\{\ell_1\ell_2,\ell_1\ell_3,\ell_2\ell_3\},
  \qquad \dim F\leq3.
\]
If \(T\neq0\), then \(T\in F^{(1)}\), so \(\dim F^{(1)}\geq1\).
Lemma~\ref{n3:lem:kernel-prolongation} gives
\[
  \rank K(T)=9\dim F-\dim F^{(1)}\leq9\cdot3-1=26.
\]
The zero term is immediate.

If the factors are linearly independent, a \(GL(V)\)-change of coordinates
reduces to \(T=x_1x_2x_3\). Then
\[
  F=\Span\{x_1x_2,x_1x_3,x_2x_3\},
  \qquad F^{(1)}=\Span(x_1x_2x_3),
\]
so the rank is \(27-1=26\).
\end{proof}

\subsection{The lower bound and a four-term decomposition}

\begin{theorem}\label{n3:thm:lower}
\[
  \ChowRank(\perm_3)\geq4.
\]
\end{theorem}

\begin{proof}
Suppose \(\perm_3=T_1+\cdots+T_r\) is a sum of \(r\) Chow terms. The
flattening is linear in the polynomial, and matrix rank is subadditive.
Corollary~\ref{n3:cor:perm-rank} and Proposition~\ref{n3:prop:term-rank}
therefore give
\[
  80=\rank K(\perm_3)
  \leq\sum_{s=1}^{r}\rank K(T_s)\leq26r.
\]
Hence
\[
  r\geq\left\lceil\frac{80}{26}\right\rceil=4.
\]
\end{proof}

Glynn's identity for \(n=3\) gives
\begin{align}\label{n3:eq:glynn}
4\perm_3={}&\prod_{j=0}^{2}(x_{0j}+x_{1j}+x_{2j})\notag\\
&-\prod_{j=0}^{2}(x_{0j}+x_{1j}-x_{2j})\notag\\
&-\prod_{j=0}^{2}(x_{0j}-x_{1j}+x_{2j})\\
&+\prod_{j=0}^{2}(x_{0j}-x_{1j}-x_{2j}).\notag
\end{align}
Each line on the right is a product of three linear forms, and its scalar and
sign may be absorbed into one factor. Thus \(\ChowRank(\perm_3)\leq4\).
Together with Theorem~\ref{n3:thm:lower}, this proves
Theorem~\ref{n3:thm:main}.

\subsection{Independent computation}

The script \texttt{perm3\_exact\_verification.py} independently constructs
the integer matrices from \eqref{n3:eq:E} and \eqref{n3:eq:delta}. Running
\begin{verbatim}
python perm3_exact_verification.py
\end{verbatim}
returns the following values.
\begin{center}
\begin{tabular}{@{}lr@{}}
\toprule
Quantity & Value\\
\midrule
\(\dim E\) & 9\\
\(\dim E^{(1)}\) & 1\\
\(\rank K(\perm_3)\) & 80\\
Koszul rank of a term with independent factors & 26\\
Flattening lower bound & 4\\
Number of terms in Glynn's decomposition & 4\\
\bottomrule
\end{tabular}
\end{center}
The program finds nonzero integer minors of orders \(80\) and \(26\) modulo
\(1000003\). The prolongation computations above give the matching
characteristic-zero upper bounds. Thus the replay is an exact cross-check,
not floating-point or probabilistic evidence. The proof of the theorem is
already complete without the computation.


The lower bound above is proved over the complex numbers. Any decomposition
over a characteristic-zero field remains a decomposition after scalar
extension to the complex numbers, while Glynn's decomposition is defined over
the prime field. The equality therefore holds over every field of
characteristic zero.

\end{document}
